\appendix
\section{Formal Derivations of Exact Structural Complexities}
\label{sec:appendix_proofs}

This appendix provides complete, self-contained algebraic derivations for the exact clause and auxiliary variable counts generated by the Boolean SAT encodings on an $N \times N$ chessboard for any integer $N \ge 2$.

\subsection{Board Geometric Decomposition and Group Cardinalities}
\label{subsec:board_decomposition}

The standard N-Queens problem on an $N \times N$ grid is partitioned into four families of cardinality constraints:
\begin{enumerate}
    \item \textbf{Row Constraints:} Exactly $N$ rows, each containing a group of $m = N$ primary variables requiring Exactly-One (1 ALO clause + 1 AMO constraint per row).
    \item \textbf{Column Constraints:} Exactly $N$ columns, each containing a group of $m = N$ primary variables requiring Exactly-One (1 ALO clause + 1 AMO constraint per column).
    \item \textbf{Main Diagonals ($i - j = d$):} The diagonal index ranges over $d \in \{-(N-1), \dots, N-1\}$. The two corner diagonals ($d = \pm (N-1)$) contain a single cell ($m=1$) and are trivially non-conflicting. For non-trivial diagonals requiring At-Most-One ($m \ge 2$):
    \begin{itemize}
        \item The main diagonal ($d = 0$) has length $m = N$.
        \item For each integer length $m \in \{2, \dots, N-1\}$, there are exactly $2$ symmetric diagonals of length $m$ (corresponding to $d = N - m$ and $d = -(N - m)$).
    \end{itemize}
    \item \textbf{Anti-Diagonals ($i + j = s$):} The anti-diagonal index ranges over $s \in \{0, \dots, 2N-2\}$. The two corner cells ($s = 0$ and $s = 2N-2$) contain $m=1$. For non-trivial anti-diagonals ($m \ge 2$):
    \begin{itemize}
        \item The central anti-diagonal ($s = N-1$) has length $m = N$.
        \item For each integer length $m \in \{2, \dots, N-1\}$, there are exactly $2$ symmetric anti-diagonals of length $m$.
    \end{itemize}
\end{enumerate}

Summing across both diagonal families, the board contains:
\begin{itemize}
    \item Exactly $2$ diagonals of length $N$.
    \item Exactly $4$ diagonals of length $m$ for each $m \in \{2, \dots, N-1\}$.
\end{itemize}
The total number of non-trivial diagonal groups requiring AMO is $2 + 4(N-2) = 4N - 6$. Combined with the $2N$ row and column groups, the entire board comprises $6N - 6$ constraint groups for $N \ge 2$.

\subsection{Exact Clause Count for Pairwise AMO Encoding}
\label{subsec:proof_pairwise}

\textbf{Theorem 1.} \textit{For the N-Queens problem with $N \ge 2$, the Pairwise (Binomial) encoding generates zero auxiliary variables and an exact total clause count of:}
\begin{equation}
C_{\text{pairwise}}(N) = 2N + N^2(N-1) + \frac{N(N-1)(2N-1)}{3}
\end{equation}

\textit{Proof.} The total clause count consists of $2N$ At-Least-One clauses (one for each row and column) plus the sum of Pairwise AMO clauses across all groups. A group of size $m$ generates $\binom{m}{2} = \frac{m(m-1)}{2}$ exclusion clauses.
\begin{enumerate}
    \item \textbf{Rows and Columns:} There are $2N$ groups of size $N$, yielding:
    \begin{equation}
    C_{\text{row,col}}^{\text{AMO}} = 2N \cdot \frac{N(N-1)}{2} = N^2(N-1)
    \end{equation}
    \item \textbf{Diagonals:} The diagonal groups comprise 2 groups of length $N$ and 4 groups of each length $m \in \{2, \dots, N-1\}$:
    \begin{align}
    C_{\text{diag}}^{\text{AMO}} &= 2 \cdot \frac{N(N-1)}{2} + \sum_{m=2}^{N-1} 4 \cdot \frac{m(m-1)}{2} \nonumber \\
    &= N(N-1) + 2 \sum_{m=2}^{N-1} (m^2 - m)
    \end{align}
    Applying the standard summation identities $\sum_{k=1}^K k = \frac{K(K+1)}{2}$ and $\sum_{k=1}^K k^2 = \frac{K(K+1)(2K+1)}{6}$ with upper bound $K = N-1$:
    \begin{align}
    \sum_{m=2}^{N-1} (m^2 - m) &= \sum_{m=1}^{N-1} (m^2 - m) \nonumber \\
    &= \frac{(N-1)N(2N-1)}{6} - \frac{(N-1)N}{2} \nonumber \\
    &= \frac{N(N-1)(2N - 4)}{6} = \frac{N(N-1)(N-2)}{3}
    \end{align}
    Substituting this back into $C_{\text{diag}}^{\text{AMO}}$:
    \begin{align}
    C_{\text{diag}}^{\text{AMO}} &= N(N-1) + \frac{2N(N-1)(N-2)}{3} \nonumber \\
    &= \frac{N(N-1)[3 + 2(N-2)]}{3} = \frac{N(N-1)(2N-1)}{3}
    \end{align}
\end{enumerate}
Summing the $2N$ ALO clauses, the row/column AMO clauses, and the diagonal AMO clauses yields:
\begin{equation}
C_{\text{pairwise}}(N) = 2N + N^2(N-1) + \frac{N(N-1)(2N-1)}{3}
\end{equation}
Expanding this polynomial yields $C_{\text{pairwise}}(N) = \frac{5}{3}N^3 - 2N^2 + \frac{7}{3}N \in \Theta(N^3)$, which proves the theorem. \hfill $\blacksquare$

\subsection{Exact Complexity for Sinz Sequential Counter AMO Encoding}
\label{subsec:proof_sequential}

\textbf{Theorem 2.} \textit{For the N-Queens problem with $N \ge 2$, the Sinz Sequential Counter encoding generates exact totals of:}
\begin{align}
A_{\text{sequential}}(N) &= 4N^2 - 6N + 2 \\
C_{\text{sequential}}(N) &= 12N^2 - 22N + 12
\end{align}
\textit{auxiliary variables and clauses, respectively.}

\textit{Proof.} In Sinz's sequential counter encoding for an At-Most-One constraint over $m \ge 2$ variables, exactly $m - 1$ auxiliary variables are instantiated, and exactly $3m - 4$ clauses are generated (comprising 1 initial implication clause, $3(m-2)$ transition and exclusion clauses, and 1 terminal exclusion clause).

\textbf{Auxiliary Variables ($A_{\text{sequential}}$):}
\begin{enumerate}
    \item \textbf{Rows and Columns:} $2N$ groups of size $N$ contribute $2N(N-1)$ auxiliary variables.
    \item \textbf{Diagonals:} The diagonal groups contribute:
    \begin{align}
    A_{\text{diag}} &= 2(N-1) + \sum_{m=2}^{N-1} 4(m-1) \nonumber \\
    &= 2(N-1) + 4 \sum_{j=1}^{N-2} j \nonumber \\
    &= 2(N-1) + 4 \cdot \frac{(N-2)(N-1)}{2} \nonumber \\
    &= 2(N-1) + 2(N-2)(N-1) = 2(N-1)^2
    \end{align}
\end{enumerate}
Summing the row/column and diagonal auxiliaries:
\begin{align}
A_{\text{sequential}}(N) &= 2N(N-1) + 2(N-1)^2 \nonumber \\
&= 2(N-1)(N + N - 1) = 2(N-1)(2N-1) \nonumber \\
&= 4N^2 - 6N + 2
\end{align}

\textbf{Clauses ($C_{\text{sequential}}$):}
The clause total includes $2N$ ALO clauses plus the AMO clauses across all groups ($3m-4$ clauses per group of size $m$):
\begin{enumerate}
    \item \textbf{Rows and Columns:} $2N(3N-4) = 6N^2 - 8N$.
    \item \textbf{Diagonals:}
    \begin{align}
    C_{\text{diag}}^{\text{AMO}} &= 2(3N-4) + \sum_{m=2}^{N-1} 4(3m-4) \nonumber \\
    &= (6N-8) + 12 \sum_{m=2}^{N-1} m - 16(N-2)
    \end{align}
    Evaluating the summation $\sum_{m=2}^{N-1} m = \frac{(N-1)N}{2} - 1 = \frac{N^2 - N - 2}{2}$:
    \begin{align}
    C_{\text{diag}}^{\text{AMO}} &= 6N - 8 + 6(N^2 - N - 2) - 16N + 32 \nonumber \\
    &= 6N^2 + (6 - 6 - 16)N + (-8 - 12 + 32) \nonumber \\
    &= 6N^2 - 16N + 12
    \end{align}
\end{enumerate}
Summing the ALO clauses, row/column AMO clauses, and diagonal AMO clauses:
\begin{align}
C_{\text{sequential}}(N) &= 2N + (6N^2 - 8N) + (6N^2 - 16N + 12) \nonumber \\
&= 12N^2 - 22N + 12
\end{align}
Both $A_{\text{sequential}}(N)$ and $C_{\text{sequential}}(N)$ are quadratic in $N$ ($\Theta(N^2)$), which proves the theorem. \hfill $\blacksquare$

\subsection{Empirical Verification of Closed-Form Expressions}
\label{subsec:verification_table}

Table~\ref{tab:appendix_verification} validates the closed-form expressions against the values produced by the benchmark model generator across representative instance sizes.

\begin{table}[htbp]
\centering
\caption{Exact theoretical vs.\ generated structural counts across instance sizes.}
\label{tab:appendix_verification}
\begin{tabular}{rrrrr}
\toprule
$N$ & $C_{\text{pairwise}}$ (Formula) & $C_{\text{pairwise}}$ (Generated) & $C_{\text{sequential}}$ (Formula) & $C_{\text{sequential}}$ (Generated) \\
\midrule
4   & 84         & 84         & 116       & 116       \\
8   & 744        & 744        & 604       & 604       \\
16  & 6,248      & 6,248      & 2,732     & 2,732     \\
20  & 12,380     & 12,380     & 4,372     & 4,372     \\
32  & 51,760     & 51,760     & 11,596    & 11,596    \\
50  & 203,770    & 203,770    & 28,912    & 28,912    \\
64  & 424,848    & 424,848    & 47,820    & 47,820    \\
100 & 1,646,900  & 1,646,900  & 117,812   & 117,812   \\
200 & 13,253,800 & 13,253,800 & 475,612   & 475,612   \\
\bottomrule
\end{tabular}
\end{table}

